\chapter{El cálculo de Feynman}
\label{ch:feynmanCalc}

En la práctica, la formulación cuantitativa de la dinámica de partículas elementales se reduce al cálculo de \emph{tasas de decaimiento} ($\decrate$) y \emph{secciones transversales de dispersión} ($\sigma$). El procedimiento involucra, en primer lugar, la evaluación de los diagramas de Feynman relevantes para determinar la \emph{amplitud invariante} ($\aminv$) para el proceso en cuestión y, posteriormente, la inserción de esta cantidad en la `\emph{regla de oro'} de Fermi para determinar $\decrate$ o $\sigma$, según corresponda. Siguiendo a Griffiths (Citation needed), en este capítulo desarrollaremos un modelo simplificado que nos permita comprender los conceptos fundamentales sin opacar los razonamientos con complicaciones algebráicas.

\section{Decaimientos y dispersión}

